\chapter{Lezione 13 --- Assegnamento degli autovalori mediante retroazione completa dello stato}

\section{Impostazione del problema}

Consideriamo un sistema lineare tempo-invariante in forma di stato
\[
\begin{cases}
\dot x(t)=Ax(t)+Bu(t),\\[4pt]
y(t)=Cx(t),
\end{cases}
\]
dove, per evitare ambiguità di notazione,
\[
x(t)\in\mathbb{R}^n,\qquad
u(t)\in\mathbb{R}^r,\qquad
y(t)\in\mathbb{R}^\ell,
\]
e
\[
A\in\mathbb{R}^{n\times n},\qquad
B\in\mathbb{R}^{n\times r},\qquad
C\in\mathbb{R}^{\ell\times n}.
\]

Vogliamo controllare il sistema attraverso una \textbf{retroazione dello stato}.

\subsection{Problema di regolazione e problema di tracking}

Negli appunti si distinguono due casi:
\begin{itemize}
    \item se il riferimento \(y^\star\) è costante, si ha un \textbf{problema di regolazione};
    \item se il riferimento è tempo-variabile,
    \[
    y^\star=y^\star(t),
    \]
    si ha un \textbf{problema di tracking}.
\end{itemize}

Nel seguito ci concentriamo soprattutto sulla parte di stabilizzazione e assegnamento degli autovalori,
che riguarda la dinamica interna del sistema retroazionato.

\section{Retroazione completa dello stato}

Supponiamo che tutto lo stato sia misurabile e disponibile.
In tal caso possiamo usare una legge di controllo del tipo
\[
u(t)=-Kx(t)+n(t),
\]
dove:
\begin{itemize}
    \item \(K\) è la matrice di guadagno di retroazione;
    \item \(n(t)\) è un ingresso esterno, usato ad esempio per il riferimento.
\end{itemize}

Nel caso SISO, \(K\) è un vettore riga:
\[
K\in\mathbb{R}^{1\times n}.
\]

\subsection*{Osservazione}

L'ipotesi cruciale è che lo stato sia \emph{tutto misurabile}.
Qui non stiamo ancora affrontando il problema dell'osservatore:
si assume direttamente di conoscere \(x(t)\).

\section{Dinamica del sistema retroazionato}

Sostituendo la legge di controllo
\[
u(t)=-Kx(t)+n(t)
\]
nell'equazione di stato,
si ottiene
\[
\dot x(t)=Ax(t)+B\bigl(-Kx(t)+n(t)\bigr).
\]

Quindi
\[
\dot x(t)=Ax(t)-BKx(t)+Bn(t),
\]
ossia
\[
\dot x(t)=(A-BK)x(t)+Bn(t).
\]

La matrice di stato del sistema retroazionato è dunque
\[
A_{\text{cl}}=A-BK.
\]

\section{Autovalori del sistema retroazionato}

Gli autovalori della matrice \(A\) sono le radici del polinomio caratteristico
\[
p_A(\lambda)=\det(\lambda I-A).
\]

Dopo la retroazione di stato, gli autovalori del sistema diventano quelli della matrice
\[
A-BK,
\]
cioè le radici di
\[
\det(\lambda I-A+BK)=0.
\]

La domanda fondamentale è quindi:

\begin{quote}
\emph{Posso progettare \(K\) in modo da assegnare arbitrariamente gli autovalori di \(A-BK\)?}
\end{quote}

\section{Teorema di assegnamento degli autovalori}

\subsection*{Caso continuo}

Negli appunti il teorema è formulato nel caso tempo continuo.

\begin{theorem}
Se la coppia \((A,B)\) è completamente controllabile, allora esiste una matrice di retroazione \(K\)
tale da assegnare arbitrariamente gli \(n\) autovalori della matrice
\[
A-BK.
\]
\end{theorem}

Nel caso continuo, controllabilità e raggiungibilità sono equivalenti,
quindi l'ipotesi può essere formulata in entrambi i modi.

\subsection*{Osservazione}

Il termine di riferimento \(n(t)\) non entra nel teorema:
esso non modifica la matrice di stato del sistema retroazionato.
Per lo studio dell'assegnamento degli autovalori si può quindi porre
\[
n(t)=0.
\]

\section{Caso SISO e forma canonica di controllo}

La dimostrazione negli appunti viene svolta nel caso SISO, cioè per
\[
\dot x=Ax+bu,
\qquad b\in\mathbb{R}^n.
\]

Se la coppia \((A,b)\) è completamente controllabile,
allora esiste una trasformazione di coordinate invertibile
\[
z=P^{-1}x
\]
che porta il sistema in \textbf{forma canonica di controllo}:
\[
\dot z=A_{cc}z+b_{cc}u.
\]

La matrice \(A_{cc}\) ha la forma
\[
A_{cc}=
\begin{bmatrix}
0 & 1 & 0 & \cdots & 0\\
0 & 0 & 1 & \cdots & 0\\
\vdots & \vdots & \vdots & \ddots & \vdots\\
0 & 0 & 0 & \cdots & 1\\
-a_0 & -a_1 & -a_2 & \cdots & -a_{n-1}
\end{bmatrix},
\]
mentre
\[
b_{cc}=
\begin{bmatrix}
0\\
0\\
\vdots\\
0\\
1
\end{bmatrix}.
\]

I coefficienti \(a_0,\dots,a_{n-1}\) sono quelli del polinomio caratteristico di \(A\):
\[
\det(\lambda I-A)
=
\lambda^n+a_{n-1}\lambda^{n-1}+\cdots+a_1\lambda+a_0.
\]

\subsection*{Osservazione sulla trasformazione di similitudine}

Poiché
\[
A_{cc}=P^{-1}AP,
\]
la trasformazione è una trasformazione di similitudine, quindi conserva gli autovalori.

In particolare,
\[
\det(\lambda I-A)=\det(\lambda I-A_{cc}).
\]

\section{Retroazione in forma canonica}

Nel nuovo sistema di coordinate si applica la retroazione
\[
u=-K_{cc}z,
\qquad
K_{cc}=
\begin{bmatrix}
k_1 & k_2 & \cdots & k_n
\end{bmatrix}.
\]

La dinamica chiusa diventa
\[
\dot z=(A_{cc}-b_{cc}K_{cc})z.
\]

Poiché
\[
b_{cc}K_{cc}
=
\begin{bmatrix}
0\\
0\\
\vdots\\
0\\
1
\end{bmatrix}
\begin{bmatrix}
k_1 & k_2 & \cdots & k_n
\end{bmatrix}
=
\begin{bmatrix}
0 & 0 & \cdots & 0\\
0 & 0 & \cdots & 0\\
\vdots & \vdots & \ddots & \vdots\\
k_1 & k_2 & \cdots & k_n
\end{bmatrix},
\]
si ottiene
\[
A_{cc}-b_{cc}K_{cc}
=
\begin{bmatrix}
0 & 1 & 0 & \cdots & 0\\
0 & 0 & 1 & \cdots & 0\\
\vdots & \vdots & \vdots & \ddots & \vdots\\
0 & 0 & 0 & \cdots & 1\\
-(a_0+k_1) & -(a_1+k_2) & \cdots & -(a_{n-2}+k_{n-1}) & -(a_{n-1}+k_n)
\end{bmatrix}.
\]

Questa è ancora una matrice in forma canonica di controllo.

\section{Polinomio caratteristico del sistema retroazionato}

Il polinomio caratteristico della matrice retroazionata è
\[
\det(\lambda I-A_{cc}+b_{cc}K_{cc})
=
\lambda^n
+(a_{n-1}+k_n)\lambda^{n-1}
+\cdots
+(a_1+k_2)\lambda
+(a_0+k_1).
\]

Supponiamo di voler assegnare gli autovalori desiderati
\[
\lambda_1,\lambda_2,\dots,\lambda_n.
\]

Allora imponiamo come polinomio desiderato
\[
(\lambda-\lambda_1)(\lambda-\lambda_2)\cdots(\lambda-\lambda_n)
=
\lambda^n+\alpha_{n-1}\lambda^{n-1}+\cdots+\alpha_1\lambda+\alpha_0.
\]

Uguagliando i coefficienti dei due polinomi, si ricava:
\[
\alpha_{n-1}=a_{n-1}+k_n,
\]
\[
\alpha_{n-2}=a_{n-2}+k_{n-1},
\]
\[
\vdots
\]
\[
\alpha_0=a_0+k_1.
\]

Quindi i guadagni si scelgono come
\[
k_1=\alpha_0-a_0,
\qquad
k_2=\alpha_1-a_1,
\qquad
\dots,
\qquad
k_n=\alpha_{n-1}-a_{n-1}.
\]

Questo mostra che, in forma canonica di controllo, i guadagni di retroazione
permettono di assegnare arbitrariamente il polinomio caratteristico.

\section{Ritorno alle coordinate originali}

Poiché
\[
z=P^{-1}x
\qquad\Longrightarrow\qquad
x=Pz,
\]
si ha
\[
u=-Kx=-KPz.
\]

Confrontando con
\[
u=-K_{cc}z,
\]
segue
\[
K_{cc}=KP.
\]

Dunque
\[
K=K_{cc}P^{-1}.
\]

Questa formula permette di tornare dal sistema in forma canonica
al sistema nelle coordinate originarie.

\section{Sufficienza dell'ipotesi di controllabilità}

La costruzione appena fatta dimostra la \textbf{sufficienza} dell'ipotesi:
se \((A,b)\) è completamente controllabile,
allora è possibile assegnare arbitrariamente tutti gli autovalori del sistema retroazionato.

\section{Necessità dell'ipotesi di controllabilità}

Negli appunti viene anche mostrato perché l'ipotesi di controllabilità è necessaria.

Supponiamo che la coppia \((A,b)\) non sia completamente controllabile
e che il sottospazio di controllabilità abbia dimensione
\[
p<n.
\]

Allora esiste un cambiamento di coordinate
\[
z=Q^{-1}x
\]
che separa la parte controllabile e quella non controllabile:
\[
z=
\begin{bmatrix}
z_c\\
z_{nc}
\end{bmatrix},
\qquad
z_c\in\mathbb{R}^p,
\qquad
z_{nc}\in\mathbb{R}^{n-p}.
\]

Nel nuovo sistema di coordinate la dinamica assume la forma
\[
\dot z
=
\begin{bmatrix}
\tilde A_c & \tilde A_{12}\\
0 & \tilde A_{nc}
\end{bmatrix}
z
+
\begin{bmatrix}
\tilde b_c\\
0
\end{bmatrix}u.
\]

Scrivendo anche la retroazione in tali coordinate,
\[
\tilde K=KQ=
\begin{bmatrix}
\tilde K_c & \tilde K_{nc}
\end{bmatrix},
\]
la matrice chiusa diventa
\[
\tilde A-\tilde b\tilde K
=
\begin{bmatrix}
\tilde A_c-\tilde b_c\tilde K_c & \tilde A_{12}-\tilde b_c\tilde K_{nc}\\
0 & \tilde A_{nc}
\end{bmatrix}.
\]

Essendo una matrice triangolare a blocchi, il polinomio caratteristico fattorizza:
\[
\det(\lambda I-\tilde A+\tilde b\tilde K)
=
\det(\lambda I_p-\tilde A_c+\tilde b_c\tilde K_c)\,
\det(\lambda I_{n-p}-\tilde A_{nc}).
\]

Quindi:
\begin{itemize}
    \item i primi \(p\) autovalori, associati alla parte controllabile, possono essere spostati;
    \item i restanti \(n-p\) autovalori, associati alla parte non controllabile, \textbf{non vengono modificati} dalla retroazione di stato.
\end{itemize}

\subsection*{Conclusione}

Se la coppia non è completamente controllabile, non posso assegnare arbitrariamente tutti gli autovalori:
posso al più assegnarne \(p\), dove \(p\) è la dimensione del sottospazio controllabile.

Questo dimostra la necessità dell'ipotesi del teorema.

\section{Interpretazione fisica}

Gli autovalori della matrice \(A\) sono determinati dalla fisica del sistema.
La matrice \(B\) descrive come l'ingresso agisce sul sistema.

La retroazione di stato modifica la dinamica attraverso il termine
\[
-BKx,
\]
quindi cambia la matrice di stato da \(A\) a
\[
A-BK.
\]

Ma questa modifica può influenzare soltanto i modi dinamici che sono effettivamente raggiungibili
dall'ingresso.

Per questo:
\begin{itemize}
    \item i modi controllabili possono essere spostati;
    \item i modi non controllabili restano dov'erano.
\end{itemize}

\section{Forma finale del problema di controllo}

Il problema di assegnamento degli autovalori si può quindi formulare così:

dato il sistema
\[
\begin{cases}
\dot x=Ax+Bu,\\
y=Cx,
\end{cases}
\]
si vuole trovare una legge di controllo
\[
u(t)=-Kx(t)+n(t)
\]
tale che la matrice chiusa
\[
A-BK
\]
abbia autovalori desiderati.

Se \((A,B)\) è completamente controllabile, questo è sempre possibile.

\section{Sintesi della lezione}

I punti chiave della lezione sono:
\begin{enumerate}
    \item si considera una retroazione completa dello stato:
    \[
    u(t)=-Kx(t)+n(t);
    \]
    \item la matrice di stato del sistema retroazionato è
    \[
    A-BK;
    \]
    \item gli autovalori del sistema chiuso sono le radici di
    \[
    \det(\lambda I-A+BK)=0;
    \]
    \item se la coppia \((A,B)\) è completamente controllabile, allora è possibile assegnare arbitrariamente tutti gli autovalori del sistema chiuso;
    \item nel caso SISO, la dimostrazione si fa portando il sistema in forma canonica di controllo;
    \item in forma canonica, i coefficienti del polinomio caratteristico dipendono linearmente dai guadagni \(k_i\);
    \item uguagliando i coefficienti del polinomio desiderato si ricavano direttamente i guadagni;
    \item tornando alle coordinate originali si ottiene il guadagno \(K\) del sistema reale;
    \item se il sistema non è completamente controllabile, soltanto gli autovalori della parte controllabile possono essere assegnati;
    \item gli autovalori non controllabili non vengono modificati dalla retroazione di stato.
\end{enumerate}